% ======================================================================
%  PART 7 -- Multivariate Time Series and Vector Autoregression (VAR)
% ======================================================================
\section{Multivariate Time Series and Vector Autoregression (VAR)}
\label{sec:multivariate}

So far each $X_t$ has been a scalar. Many real systems are
$d$-dimensional: a basket of normalised stock prices (Apple, Nasdaq-100, S\&P
500), or the longitude/latitude/temperature of an ocean drifter. We now let
$\{X_t\}$ take values in $\Reals^d$ and ask two questions: (i) how do we
\emph{describe} the second-order structure --- now a matrix-valued
autocovariance and a matrix-valued spectrum, with genuine \textbf{cross}-terms
between components, and (ii) how do we \emph{model} it. The canonical model is
the vector autoregression VAR($p$), and the key trick is that
\emph{every VAR($p$) is a VAR(1)} in disguise (companion form), so its theory
reduces to that of a single matrix $\Phi_1$ (or $F$).

% ----------------------------------------------------------------------
\subsection{Definitions}

\begin{definition}{Multivariate time series}{p7-mvts}
A \textbf{multivariate time series} is a real $d$-vector-valued discrete-time
stochastic process
\[
X_t=\big(X_t^{(1)},X_t^{(2)},\dots,X_t^{(d)}\big)\Tr,\qquad t\in\Ints,
\]
whose marginal processes $\{X_t^{(1)}\},\dots,\{X_t^{(d)}\}$ are each univariate
time series. When $d=2$ we call it \textbf{bivariate}.
\end{definition}

\begin{definition}{Multivariate second-order stationarity}{p7-mvstat}
$\{X_t\}$ is \textbf{multivariate second-order (weakly) stationary} if for all
$t,s,\tau\in\Ints$,
\[
\E[X_t]=\E[X_s],\qquad
\Cov(X_{t+\tau},X_t)=\Cov(X_{s+\tau},X_s),\qquad
\trace\big(\Var(X_t)\big)<\infty .
\]
Equivalently: each marginal $\{X_t^{(j)}\}$ is second-order stationary
\textbf{and} every pair has a shift-invariant cross-covariance,
$\Cov\!\big(X_{t+\tau}^{(j)},X_t^{(k)}\big)=\Cov\!\big(X_{s+\tau}^{(j)},X_s^{(k)}\big)$
for all $1\le j,k\le d$. (Here $\Cov(U,V)$ of two vectors is the matrix with
$(j,k)$ entry $\Cov(U^{(j)},V^{(k)})$.)
\end{definition}

\begin{definition}{Matrix autocovariance sequence (ACVS)}{p7-acvs}
For a multivariate second-order stationary $\{X_t\}$, the
\textbf{(matrix) autocovariance sequence} is the $d\times d$ matrix-valued
sequence
\[
\Gamma_\tau=\Cov(X_{t+\tau},X_t),\qquad \tau\in\Ints ,
\]
with $(j,k)$ entry $\acvs_\tau^{(j,k)}=\Cov\!\big(X_{t+\tau}^{(j)},X_t^{(k)}\big)$.
\begin{itemize}
  \item \textbf{The lag $\tau$ sits in the FIRST argument.} In the scalar case
        this is immaterial, but here it determines the convention for
        $\Gamma_{-\tau}=\Gamma_\tau\Tr$ below.
  \item When $j=k$, $\acvs_\tau^{(j,j)}$ is the ordinary ACVS of the marginal
        $\{X_t^{(j)}\}$.
  \item When $j\ne k$, $\acvs_\tau^{(j,k)}$ is called the
        \textbf{cross-covariance sequence} (CCS) between components $j$ and $k$.
\end{itemize}
\end{definition}

\begin{definition}{Cross-correlation sequence (CCS)}{p7-ccs}
The \textbf{cross-correlation sequence} between components $j$ and $k$ is
\[
\acf_\tau^{(j,k)}
=\frac{\acvs_\tau^{(j,k)}}{\sqrt{\acvs_0^{(j,j)}\,\acvs_0^{(k,k)}}} .
\]
For $j=k$ this is the ordinary ACF (so $\acf_0^{(j,j)}=1$). For $j\ne k$ it need
\emph{not} equal $1$ at $\tau=0$ and need \emph{not} be symmetric in $\tau$; it
satisfies $\acf_\tau^{(j,k)}=\acf_{-\tau}^{(k,j)}$.
\end{definition}

\begin{definition}{Multivariate white noise $\mathrm{WN}(0,\Sigma)$}{p7-wn}
A $d$-variate series $\{\eps_t\}$ is \textbf{white noise with mean $0$ and
covariance matrix $\Sigma$}, written $\{\eps_t\}\sim\mathrm{WN}(0,\Sigma)$, iff
it is stationary with mean vector $0$ and
\[
\Gamma_\tau^{(\eps)}=
\begin{cases}\Sigma & \tau=0,\\[2pt] 0 & \tau\ne0.\end{cases}
\]
\begin{itemize}
  \item There \emph{can} be correlation between different components at lag $0$
        (off-diagonal entries of $\Sigma$): the components are uncorrelated
        \textbf{across time} but possibly correlated \textbf{contemporaneously}.
  \item Often one further assumes $\Sigma$ \emph{diagonal} (componentwise
        independent noise), but this is an extra assumption, not part of the
        definition.
\end{itemize}
\end{definition}

\begin{definition}{Uncorrelated processes}{p7-uncorr}
Two jointly second-order stationary processes $\{X_t^{(1)}\}$ and
$\{X_t^{(2)}\}$ are \textbf{uncorrelated} if their cross-covariance vanishes at
\emph{every} lag:
\[
\acvs_\tau^{(1,2)}=0\qquad\text{for all }\tau\in\Ints .
\]
\end{definition}

\begin{definition}{Spectral density matrix function}{p7-sdm}
Assume $\{\acvs_\tau^{(j,k)}\}_{\tau\in\Ints}\in\ell^1$ for all $1\le j,k\le d$
(absolutely summable cross-covariances). The \textbf{spectral density matrix}
is the $d\times d$ matrix-valued function
\[
S(f)=\sum_{\tau\in\Ints}\Gamma_\tau\,e^{-2\pi i\tau f},
\qquad f\in[-\tfrac12,\tfrac12],
\]
with inverse $\displaystyle \Gamma_\tau=\int_{-1/2}^{1/2}S(f)\,e^{2\pi if\tau}\,df$.
\begin{itemize}
  \item The $(j,k)$ entry is $S_{j,k}(f)$. For $j=k$ it is the ordinary
        (real-valued) spectral density $\sdf_{j}(f)$.
  \item For $j\ne k$, $S_{j,k}(f)$ is the \textbf{cross-spectral density}, which
        is \emph{complex-valued} in general.
\end{itemize}
\end{definition}

\begin{definition}{Coherence}{p7-coh}
The \textbf{coherence} between components $j$ and $k$ is
\[
r_{j,k}(f)=\frac{S_{j,k}(f)}{\sqrt{S_{j,j}(f)\,S_{k,k}(f)}}\in\Comp ,
\]
the frequency-domain analogue of a correlation between the increments $dZ_j$ and
$dZ_k$. One reports:
\begin{itemize}
  \item its \textbf{magnitude} $\abs{r_{j,k}(f)}$ (or magnitude-squared
        $\abs{r_{j,k}(f)}^2\in[0,1]$), the strength of frequency-by-frequency
        association;
  \item its \textbf{phase} $\arg r_{j,k}(f)$, the ``sign'' / lead--lag angle in
        the complex plane.
\end{itemize}
\end{definition}

\begin{definition}{Vector autoregression VAR($p$)}{p7-var}
Let $\{\eps_t\}\sim\mathrm{WN}(0,\Sigma)$ be $d$-variate. A process $\{X_t\}$ is
a \textbf{vector autoregressive process of order $p$}, VAR($p$), if
\[
X_t=\Phi_1 X_{t-1}+\Phi_2 X_{t-2}+\cdots+\Phi_p X_{t-p}+\eps_t ,
\qquad \Phi_j\in\Reals^{d\times d},\ \Phi_p\ne0 .
\]
Defining the matrix polynomial $\Phi(z)=I_d-\sum_{j=1}^p\Phi_j z^j$, this is
written compactly as
\[
\Phi(\Bsh)X_t=\eps_t ,
\]
where $\Bsh$ is the backshift operator. The same regressors (past values of
the \emph{whole} vector $X_t$) appear in every component equation; the
off-diagonal entries of the $\Phi_j$ are exactly the cross-component dynamics.
\end{definition}

\begin{definition}{Companion form of a VAR($p$)}{p7-companion}
Stack the last $p$ vectors into $Y_t=(X_t\Tr,X_{t-1}\Tr,\dots,X_{t-p+1}\Tr)\Tr
\in\Reals^{dp}$. Then a VAR($p$) is equivalently the VAR(1)
\[
\underbrace{Y_t}_{dp\times1}
=\underbrace{F}_{dp\times dp}\,\underbrace{Y_{t-1}}_{dp\times1}
+\underbrace{U_t}_{dp\times1},
\qquad
F=\begin{pmatrix}
\Phi_1 & \Phi_2 & \cdots & \Phi_{p-1} & \Phi_p\\
I_d & 0 & \cdots & 0 & 0\\
0 & I_d & & \vdots & \vdots\\
\vdots & & \ddots & 0 & 0\\
0 & \cdots & 0 & I_d & 0
\end{pmatrix},
\quad
U_t=\begin{pmatrix}\eps_t\\0\\\vdots\\0\end{pmatrix}.
\]
The top block row reproduces the VAR($p$) equation; the lower rows are
identities $X_{t-i}=X_{t-i}$ that shift the stack. Recover the original process
via $X_t=G Y_t$ with $G=(I_d\ \ 0\ \cdots\ 0)\in\Reals^{d\times dp}$.
\end{definition}

\begin{definition}{Stability of a VAR($p$)}{p7-stable}
A VAR($p$) is \textbf{stable} if the roots of the \textbf{reverse
characteristic polynomial}
\[
\det\!\big(I_d-\Phi_1 z-\Phi_2 z^2-\cdots-\Phi_p z^p\big)=0
\]
all lie \emph{strictly outside} the complex unit circle ($\abs{z}>1$).
Equivalently, the companion matrix $F$ has all eigenvalues of modulus
\emph{strictly less than one} ($\abs{\lambda}<1$). Stability $\Rightarrow$
stationarity (but not conversely in general).
\end{definition}

% ----------------------------------------------------------------------
\subsection{Theorems \& Propositions}

\begin{proposition}{Symmetry of the matrix ACVS}{p7-gammasym}
For a multivariate second-order stationary process,
\[
\acvs_\tau^{(j,k)}=\acvs_{-\tau}^{(k,j)}
\qquad\Longleftrightarrow\qquad
\Gamma_{-\tau}=\Gamma_\tau\Tr .
\]
\textit{Proof idea:} $\acvs_\tau^{(j,k)}=\Cov(X_{t+\tau}^{(j)},X_t^{(k)})
=\Cov(X_t^{(k)},X_{t+\tau}^{(j)})$ (covariance is symmetric in its arguments)
$=\Cov(X_{s-\tau}^{(k)},X_s^{(j)})$ (stationarity, set $s=t+\tau$)
$=\acvs_{-\tau}^{(k,j)}$.\\
\textit{Why:} the cross-covariance need \emph{not} be even in $\tau$
($\acvs_\tau^{(j,k)}\ne\acvs_{-\tau}^{(j,k)}$ in general); only the
\emph{transpose} symmetry $\Gamma_{-\tau}=\Gamma_\tau\Tr$ holds. This asymmetry
is precisely what encodes lead--lag relationships between channels.
\end{proposition}

\begin{proposition}{Positive semi-definiteness of the matrix ACVS}{p7-psd}
The matrix sequence $\{\Gamma_\tau\}$ is positive semi-definite: for every
$n\in\Nat$ and every $a_1,\dots,a_n\in\Reals^d$,
\[
\sum_{j=1}^n\sum_{k=1}^n a_j\Tr\,\Gamma_{k-j}\,a_k\ge0 .
\]
\textit{Proof idea:} the left side equals $\Var\!\big(\sum_{j}a_j\Tr X_{t_j}\big)\ge0$
for suitable times (the scalar random variable $\sum_j a_j\Tr X_{t_j}$ has
non-negative variance).\\
\textit{Why:} this is the defining algebraic constraint a valid matrix ACVS must
satisfy; a candidate matrix sequence violating it cannot be the autocovariance
of any process.
\end{proposition}

\begin{proposition}{Contemporaneous correlation induced by shared noise}{p7-contemp}
Let $X_t=(X_t^{(1)},X_t^{(2)})\Tr$ be second-order stationary with ACVS
$\Gamma_\tau^{(X)}$, and let $\{\phi_t\}$ be a univariate mean-zero white noise
(variance $\sigma_\phi^2$) independent of $\{X_t\}$. Define
$Y_t=X_t+\phi_t\mathbf{1}$ (the \emph{same} scalar noise added to both
components), where $\mathbf{1}=(1,1)\Tr$. Then
\[
\Gamma_\tau^{(Y)}=\Gamma_\tau^{(X)}+\sigma_\phi^2\,\delta_{\tau,0}\,\mathbf 1\mathbf 1\Tr ,
\qquad
\mathbf 1\mathbf 1\Tr=\begin{pmatrix}1&1\\1&1\end{pmatrix}.
\]
\textit{Proof idea:} $\Gamma_\tau^{(Y)}=\Cov(X_{t+\tau}+\phi_{t+\tau}\mathbf1,
X_t+\phi_t\mathbf1)$; independence kills the $X$--$\phi$ cross terms, and
$\Cov(\phi_{t+\tau},\phi_t)=\sigma_\phi^2\delta_{\tau,0}$ contributes
$\sigma_\phi^2\delta_{\tau,0}$ to \emph{every} entry.\\
\textit{Why:} adding a common shock creates lag-$0$ correlation between
$Y^{(1)}$ and $Y^{(2)}$ \emph{regardless} of the correlation in $X$. This is the
prototypical mechanism behind ``everything moves together'' contemporaneous
correlation.
\end{proposition}

\begin{proposition}{Hermitian symmetry of the spectral matrix}{p7-herm}
For all $f\in[-\tfrac12,\tfrac12]$ and all $1\le j,k\le d$,
\[
S_{j,k}(f)=\conj{S_{k,j}(f)},\qquad
S_{j,k}(f)=\conj{S_{j,k}(-f)},
\]
equivalently in matrix form
\[
S(f)=S(f)\Herm,\qquad S(f)=S(-f)\Tr ,
\]
and $S(f)$ is positive semi-definite for every $f$.
\textit{Proof idea:} substitute $\Gamma_{-\tau}=\Gamma_\tau\Tr$ into
$S(f)=\sum_\tau\Gamma_\tau e^{-2\pi i\tau f}$ and relabel $\tau\to-\tau$; the
conjugate of $e^{-2\pi i\tau f}$ flips the sign in the exponent.\\
\textit{Why:} $S(f)\Herm=S(f)$ (Hermitian) guarantees real diagonal spectra and
makes coherence well defined; positive semi-definiteness is the frequency-domain
counterpart of Proposition~\ref{prop:p7-psd}.
\end{proposition}

\begin{theorem}{Multivariate spectral representation theorem}{p7-spectral}
Let $\{X_t\}$ be a vector-valued discrete-time stationary process with mean
$\mu$. There exists a vector-valued orthogonal-increment process $\{Z(f)\}$ on
$[-\tfrac12,\tfrac12]$ such that
\[
X_t=\mu+\int_{-1/2}^{1/2} e^{2\pi i f t}\,dZ(f).
\]
If moreover $\{\acvs_\tau^{(j,k)}\}\in\ell^1$, then for $f,f'\in[-\tfrac12,\tfrac12]$,
\[
\E[dZ(f)]=0,\qquad
\Var\big(dZ(f)\big)=S(f)\,df,\qquad
\Cov\big(dZ(f),dZ(f')\big)=0\ \ (f\ne f'),
\]
where $S(f)$ is the spectral density \emph{matrix} and the $0$ in the last
identity is the zero matrix.\\
\textit{Proof idea:} the multivariate analogue of the scalar Cram\'er/spectral
representation; the increment process $Z(f)$ carries the random ``Fourier
content'' of each component, with covariance structure given by the spectral
matrix.\\
\textit{Why:} it justifies thinking of a multivariate stationary process as a
superposition of uncorrelated random sinusoids whose (matrix) variance is
$S(f)\,df$; coherence is then literally the correlation of the increments
$dZ_j(f),dZ_k(f)$.
\end{theorem}

\begin{proposition}{Cross-spectrum of a pure delay process}{p7-delayspec}
Let $\{Y_t\}$ be stationary and $X_t=(Y_t,\,Y_{t-\nu})\Tr$ for fixed
$\nu\in\Ints$. Then
\[
S_{1,1}(f)=S_{2,2}(f)=\sdf_Y(f),\qquad
S_{1,2}(f)=\sdf_Y(f)\,e^{2\pi i f\nu},
\]
so the coherence is $r_{1,2}(f)=e^{2\pi i f\nu}$, of magnitude $1$ and phase
$2\pi\nu f$.
\textit{Proof idea:} $\acvs_\tau^{(1,2)}=\Cov(Y_{t+\tau},Y_{t-\nu})
=\acvs_{\tau+\nu}^{(Y)}$; Fourier-transforming and shifting the summation index
pulls out $e^{2\pi i f\nu}$.\\
\textit{Why:} a pure time shift between two channels produces \emph{perfect}
coherence (magnitude $1$) with a phase that is \emph{linear in $f$} with slope
$2\pi\nu$ --- the textbook signature of a lead/lag of $\nu$ samples.
\end{proposition}

\begin{theorem}{Everything is a VAR(1) (companion reduction)}{p7-everyvar1}
Any VAR($p$) process $\{X_t\}$ admits the VAR(1) representation
$Y_t=FY_{t-1}+U_t$ of Definition~\ref{def:p7-companion}, with
$\{U_t\}$ a $dp$-variate white noise whose only nonzero covariance block is the
top-left $\Sigma$. Hence all VAR($p$) properties (stability, MA($\infty$)
expansion, covariances) follow from the VAR(1) theory applied to $F$ and then
projected back by $X_t=GY_t$.\\
\textit{Proof idea:} write out the stacked vector $Y_t$; its top block equals
the VAR($p$) recursion and its remaining blocks are trivial identities, giving
exactly $F$ and $U_t$.\\
\textit{Why:} it lets us prove one clean VAR(1) result and reuse it for every
order $p$, e.g.\ the stability test becomes ``eigenvalues of $F$ inside the unit
circle''.
\end{theorem}

\begin{proposition}{MA($\infty$) representation of a stable VAR(1)}{p7-mainf}
If the VAR(1) $X_t=\Phi_1 X_{t-1}+\eps_t$ is stable, then
\[
X_t=\sum_{j=0}^{\infty}\Phi_1^{\,j}\,\eps_{t-j}.
\]
\textit{Proof idea:} back-substitute repeatedly,
$X_t=\Phi_1(\Phi_1 X_{t-2}+\eps_{t-1})+\eps_t=\Phi_1^2X_{t-2}+\Phi_1\eps_{t-1}+\eps_t=\cdots$;
stability ($\Phi_1^j\to0$) makes the remainder term vanish and the series
converge (in mean square).\\
\textit{Why:} the matrix geometric series is the multivariate causal linear
filter of the noise --- it underlies the covariance formula below and all
forecasting.
\end{proposition}

\begin{proposition}{Covariance of a stable VAR(1)}{p7-var1cov}
For a stable VAR(1) with $\{\eps_t\}\sim\mathrm{WN}(0,\Sigma)$ and $\tau\ge0$,
\[
\Gamma_\tau=\Cov(X_{t+\tau},X_t)
=\sum_{k=0}^{\infty}\Phi_1^{\,k+\tau}\,\Sigma\,\big(\Phi_1^{\,k}\big)\Tr .
\]
\textit{Proof idea:} substitute the MA($\infty$) form into
$\Cov\big(\sum_j\Phi_1^{j}\eps_{t+\tau-j},\sum_k\Phi_1^{k}\eps_{t-k}\big)$;
white noise makes only the terms with $t+\tau-j=t-k$, i.e.\ $j=k+\tau$, survive,
each contributing $\Phi_1^{k+\tau}\Sigma(\Phi_1^{k})\Tr$.\\
\textit{Why:} closed-form second-order structure of a VAR(1); for $\tau=0$ it
gives $\Gamma_0=\sum_k\Phi_1^k\Sigma(\Phi_1^k)\Tr$, the stationary covariance,
which also solves the discrete Lyapunov equation
$\Gamma_0=\Phi_1\Gamma_0\Phi_1\Tr+\Sigma$.
\end{proposition}

\begin{corollary}{Covariance of a VAR($p$) via the companion form}{p7-varpcov}
For a stable VAR($p$) with companion matrix $F$ and
$\Sigma_U=\Var(U_t)$ (top-left block $\Sigma$, all else $0$),
\[
\Gamma_\tau^{(X)}=\Cov(X_{t+\tau},X_t)
=G\Big(\sum_{k=0}^{\infty}F^{\,k+\tau}\,\Sigma_U\,\big(F^{\,k}\big)\Tr\Big)G\Tr,
\qquad G=(I_d\ 0\ \cdots\ 0).
\]
\textit{Proof idea:} apply Proposition~\ref{prop:p7-var1cov} to the VAR(1)
$Y_t=FY_{t-1}+U_t$ to get $\Gamma_\tau^{(Y)}$, then use $X_t=GY_t$ so
$\Gamma_\tau^{(X)}=G\Gamma_\tau^{(Y)}G\Tr$.\\
\textit{Why:} reduces the order-$p$ covariance to one VAR(1) computation plus a
projection.
\end{corollary}

\begin{theorem}{Yule--Walker equations for a VAR($p$)}{p7-yw}
For a stable VAR($p$) and $\tau\ge0$,
\[
\Gamma_\tau=\Phi_1\Gamma_{\tau-1}+\Phi_2\Gamma_{\tau-2}
+\cdots+\Phi_p\Gamma_{\tau-p}+\delta_{\tau,0}\,\Sigma .
\]
In particular for VAR(1): $\Gamma_\tau=\Phi_1\Gamma_{\tau-1}+\delta_{\tau,0}\Sigma$.
\textit{Proof idea:} multiply the VAR equation on the right by $X_t\Tr$ and take
expectations: $\Cov(X_{t+\tau},X_t)=\Cov(\sum_i\Phi_i X_{t+\tau-i}+\eps_{t+\tau},X_t)$;
the noise term $\Cov(\eps_{t+\tau},X_t)$ contributes $\Sigma$ only when
$\tau=0$ (since $X_t$ depends only on $\eps_s$ for $s\le t$).\\
\textit{Why:} given $\Sigma$ and the $\Phi_j$, solving the first $p+1$ equations
yields $\Gamma_0,\dots,\Gamma_p$, and the recursion then generates all
$\Gamma_\tau$ for $\tau>p$. Run in reverse, it estimates the $\Phi_j$ from sample
covariances (VAR fitting).
\end{theorem}

% ----------------------------------------------------------------------
% ----------------------------------------------------------------------
